\documentclass[UTF8,a4paper,zihao=-4,fontset=fandol]{ctexart}
\usepackage[a4paper,left=23mm,right=27mm,top=20mm,bottom=21mm]{geometry}
\usepackage{amsmath,amssymb,booktabs,tabularx,array}
\usepackage{graphicx,xcolor,tikz}
\usetikzlibrary{arrows.meta,positioning,calc}
\usepackage{caption,float}
\usepackage[hidelinks,unicode]{hyperref}
\usepackage{url}
\definecolor{ink}{HTML}{17364B}
\definecolor{accent}{HTML}{236B81}
\definecolor{pale}{HTML}{EEF4F6}
\definecolor{muted}{HTML}{50616A}
\ctexset{section={format=\Large\bfseries\color{ink},beforeskip=1ex,afterskip=1ex},subsection={format=\large\bfseries,beforeskip=1ex,afterskip=.7ex}}
\captionsetup{font=small,labelfont=bf,labelsep=quad,skip=5pt}
\setlength{\parindent}{2em}
\setlength{\parskip}{.36em}
\setlength{\emergencystretch}{2em}
\setlength{\tabcolsep}{4pt}
\renewcommand{\arraystretch}{1.16}
\linespread{1.17}
\pagestyle{plain}
\graphicspath{{../assets/figures/}}
\hypersetup{pdfauthor={},pdfsubject={DRAM单粒子效应仿真研究},pdfcreator={XeLaTeX}}
\newcommand{\um}{\,\mu\mathrm{m}}
\newcommand{\ns}{\,\mathrm{ns}}
\newcommand{\mV}{\,\mathrm{mV}}
\newcommand{\fC}{\,\mathrm{fC}}
\newcommand{\V}{\,\mathrm{V}}
\newcommand{\fig}[4]{\begin{figure}[H]\centering\includegraphics[width=#2\linewidth]{#1}\caption{#3}\label{#4}\end{figure}}
\newcommand{\atable}[3]{\begin{table}[H]\centering\small\caption{#2}\label{#3}\input{../assets/tables/#1.tex}\end{table}}
\newcommand{\pageheading}[1]{\clearpage\section{#1}}
\newcommand{\smallnote}[1]{\par\noindent{\small\color{muted}#1}\par}
\newcommand{\modeldiagram}{%
\begin{tikzpicture}[font=\small,box/.style={draw=accent,fill=pale,rounded corners=1pt,minimum height=1.1cm,text width=3.35cm,align=center},arr/.style={-{Latex[length=2mm]},draw=accent,thick}]
\node[box] (bl) {位线 BL\\$C_b=100\,\mathrm{fF}$};
\node[box,right=1.0cm of bl] (mos) {二维访问晶体管\\栅长 $L$};
\node[box,right=1.0cm of mos] (sn) {存储节点 SN\\$C_s=10\,\mathrm{fF}$};
\node[above=.6cm of mos] (wl) {字线 WL：控制导通};
\draw[<->,thick,draw=accent] (bl)--(mos);
\draw[<->,thick,draw=accent] (mos)--(sn);
\draw[arr] (wl)--(mos);
\node[below=.65cm of sn,align=center] {初态 $0.1/0.9\,\mathrm V$\\分别对应存储0/1};
\node[below=.65cm of bl,align=center] {读出前预充至$0.5\,\mathrm V$\\经$50\,\Omega$支路};
\end{tikzpicture}}
\newcommand{\workflowdiagram}{%
\begin{tikzpicture}[font=\small,box/.style={draw=accent,fill=pale,minimum height=1.05cm,text width=3.55cm,align=center},arr/.style={-{Latex[length=2mm]},draw=accent,thick},node distance=.7cm and .7cm]
\node[box] (request) {中文参数请求\\检查范围和单位};
\node[box,right=of request] (g4) {Geant4\\逐步电离沉积};
\node[box,right=of g4] (source) {等效电荷源\\总量与径迹};
\node[box,below=of source] (tcad) {TCAD瞬态\\信号、对照和参考};
\node[box,left=of tcad] (metrics) {电荷变化与读出裕量\\$\Delta Q,\ M,\ \Delta M$};
\node[box,left=of metrics] (score) {固定预测评分\\留出与外推};
\draw[arr] (request)--(g4);\draw[arr] (g4)--(source);\draw[arr] (source)--(tcad);\draw[arr] (tcad)--(metrics);\draw[arr] (metrics)--(score);
\end{tikzpicture}}
\newcommand{\paramtable}{%
\begin{tabularx}{\linewidth}{@{}p{3.6cm}X@{}}\toprule
参数&设定\\\midrule
栅长、栅氧厚度&$L=0.18/0.25/0.35/0.45\um$；主矩阵$5\,\mathrm{nm}$\\
硅体尺寸&宽$L+0.6\um$，深$1.2\um$，有效第三维厚度$1\um$\\
体区、源漏掺杂&P型$10^{17}\,\mathrm{cm}^{-3}$；源漏高斯峰值$10^{20}\,\mathrm{cm}^{-3}$\\
源漏深度参数&$0.12\um$处指定$10^{17}\,\mathrm{cm}^{-3}$，横向因子0.4\\
栅功函数、温度&$4.6\,\mathrm{eV}$；$300\,\mathrm K$\\
电路&存储电容$10\,\mathrm{fF}$；位线与参考电容$100\,\mathrm{fF}$\\
初态&信号SN为$0.1/0.9\V$；中性参考SN为$0.5\V$\\
预充、读出&预充通路$50\,\Omega$，关断$10^{20}\,\Omega$；字线终值$1.8\V$\\\bottomrule
\end{tabularx}}
\newcommand{\timingtable}{%
\begin{tabularx}{\linewidth}{@{}lXlX@{}}\toprule
时间/ns&操作&时间/ns&操作\\\midrule
0.6&释放初态钳位&2.8&提取保持段电荷\\
0.95&检查初态电压&2.8--2.9&位线预充\\
1.0&粒子入射&3.0&断开预充\\
&&3.05--3.10&字线上升，4.0读出\\\bottomrule
\end{tabularx}}
\newcommand{\metricdefs}{%
\begin{align}
\Delta Q&=C_s\,[V_{{\rm SN},e}(2.8\ns)-V_{{\rm SN},c}(2.8\ns)],\\
D_e&=V_{{\rm BL},e}(4\ns)-V_{{\rm BL},n}(4\ns),\qquad M_e=(2b-1)D_e,\\
\Delta M&=M_c-M_e.
\end{align}}
\newcommand{\scoretable}{\begin{tabular}{lllr rrr}
\toprule
集合 & 位 & 目标 & 基准RMSE & 岭回归MAE & RMSE & 最大\\
\midrule
留出 & 0 & $\Delta Q$/fC & 1.9560 & 0.0566 & 0.0883 & 0.1816\\
留出 & 0 & $\Delta M$/mV & 18.4480 & 0.5270 & 0.7984 & 1.6456\\
留出 & 1 & $\Delta Q$/fC & 2.1549 & 0.0618 & 0.0946 & 0.1989\\
留出 & 1 & $\Delta M$/mV & 20.0666 & 0.5498 & 0.8380 & 1.7654\\
外推 & 0 & $\Delta Q$/fC & 2.0254 & 0.0528 & 0.0680 & 0.1364\\
外推 & 0 & $\Delta M$/mV & 18.8677 & 0.4050 & 0.5194 & 1.0753\\
外推 & 1 & $\Delta Q$/fC & 2.1983 & 0.0517 & 0.0672 & 0.1442\\
外推 & 1 & $\Delta M$/mV & 20.2199 & 0.3877 & 0.5049 & 1.1051\\
\bottomrule\end{tabular}
}

\hypersetup{pdftitle={DRAM单粒子读出响应与冻结预测验证}}
\begin{document}
\begin{center}
{\LARGE\bfseries\color{ink}DRAM单粒子读出响应与冻结预测验证\par}
\vspace{.5em}{指导教师：李洋\qquad 2026年10月11日}
\end{center}
\subsection*{摘要}
采用Geant4与TCAD计算简化1T1C单元的粒子响应。四种栅长、两类粒子和两个方向共64个独立事件，按存储0/1得到128个配对响应。5 mV判据下，96项正确、20项判据误读、8项临界、4项数值不确定。四个岭回归模型在测试前固定，独立测试的电荷RMSE为0.067155至0.094644 fC，裕量损失RMSE为0.504929至0.838043 mV。氧化层变体的响应差异较小；边界细化有一项分类跨界，表明连续量稳定不足以确定阈值附近类别。结果适用于本次未校准二维模型和固定入射条件。
\smallnote{关键词：DRAM；单粒子效应；Geant4；TCAD；读出裕量；岭回归}

\section{引言}
DRAM读出依靠存储电容与位线的电荷共享\cite{keeth2001}。粒子在硅中产生载流子后，漂移、扩散、复合及端口收集共同改变节点电压\cite{label1996}。生成电荷总量与存储节点净变化不同，需要同时观察粒子增量和读出差分。

陈善强、师立勤用TCAD与Geant4评估SRAM单粒子效应\cite{chen2011}；Xu等研究了用于电路仿真的单粒子瞬态回归模型\cite{xu2021}。不同激励方式的二极管瞬态比较还表明，电荷产生的空间分布会影响响应\cite{ryder2021}。本研究采用总量等效源，在统一电学条件下比较尺寸和位值，并用未参加拟合的尺寸检验固定预测。

\begin{figure}[H]\centering\workflowdiagram\caption{计算与评价流程。}\end{figure}
\atable{progress}{主矩阵分组。每个栅长包含两粒子、两方向及各4个事件。}{p:matrix}

\pageheading{器件与粒子源}
\begin{table}[H]\centering\small\caption{二维1T1C模型参数。各数值为探索设定。}\paramtable\end{table}
器件采用漂移扩散模型，联立泊松、电子与空穴连续性方程和电路方程，计入Fermi统计、OldSlotboom带隙修正、掺杂和电场相关迁移率、SRH与Auger复合；瞬态使用后向欧拉离散。源漏深度参数不等于净掺杂结深，四个栅长也不表示商业工艺代际。

Geant4 10.7.3使用电磁option4及步长限制\cite{g4em}。$5\,\mathrm{MeV}$的$\alpha$粒子或质子从$(L+0.45,0,0)\um$入射，方向为$(-\sin\theta,\cos\theta,0)$，$\theta=0^\circ/30^\circ$。材料为硅和外部真空；有限直径迹长度为
\begin{equation}\ell=\begin{cases}1.2\um,&\theta=0^\circ,\\\min\!\left(1.2/\cos\theta,(L+0.45)/\sin\theta\right)\um,&\theta=30^\circ.\end{cases}\end{equation}
其中$L$以$\mu\mathrm m$计。$0.18\um$、$30^\circ$的径迹从侧边退出，长度$1.26\um$。
\begin{equation}E_{\rm ion}=\sum_j\max(0,E_{{\rm dep},j}-E_{{\rm NIEL},j}),\qquad Q_{\rm dep}=q\frac{10^6E_{\rm ion}[\mathrm{MeV}]}{W[\mathrm{eV}]},\quad W=3.6\,\mathrm{eV}.\end{equation}
取$q=1.602176634\times10^{-19}\,\mathrm C$。TCAD将总量沿直径迹均匀分配，横向高斯尺度$0.03\um$，线电荷密度$Q_{\rm dep}/\ell$以$\mathrm{pC}/\mu\mathrm m$传入。源积分须计入$1\um$有效厚度，并满足$|Q_{\rm src}-Q_{\rm dep}|/Q_{\rm dep}\leq5\%$；零沉积单独检查。该检查核对注入总量，不恢复完整三维沉积分布。
\begin{table}[H]\centering\small\caption{保持与读出时序。}\timingtable\end{table}

\pageheading{读出评价与响应曲线}
用$e,c,n$分别表示事件、同存储位无粒子对照和中性参考，$b\in\{0,1\}$表示存储位。
\metricdefs
$M_c$由对照BL和相同中性参考计算。$\Delta Q$表示存储电容相对对照的变化，$\Delta M>0$表示裕量下降。中性参考决定绝对读出差分，同位对照提取粒子增量。

初态误差须不超过$1\mV$，预充误差不超过$10\,\mu\mathrm V$；检查还覆盖$4\ns$终点、有限数据、KCL、四端口闭合及电荷积分。匹配参考先与双支路模型比较两侧BL及其差分。历史理想参考留下约12.37 mV偏移；预充电阻由100 Ω改为50 Ω后，独立检查的误差从$-13.297362$变为$-4.731204\,\mu\mathrm V$，据此确定统一主模型。
\fig{readout-l018.png}{1}{$0.18\um$下两种位值的SN与BL响应。粒子、角度、同位对照和中性参考分别标示。}{p:curves}
在四栅长范围内，$\alpha$存储1的裕量为$-13.213$至$0.311\mV$，质子存储1为$31.781$至$35.651\mV$。$\alpha$存储0的$\Delta M$为负，表示相对于无粒子对照的裕量增加。逐分组范围见附表数据。

\pageheading{分类结果}
判据采用$M>\epsilon$为正确、$M<-\epsilon$为误读、$|M|\leq\epsilon$为临界。主阈值$\epsilon=5\mV$，同时比较$2.5/10\mV$。接近边界但未完成规定细化，或细化后分类改变的响应，保留为数值不确定。
\atable{thresholds}{完整128响应在三阈值下的状态。}{p:thresholds}
\fig{stratum-counts.png}{.95}{5 mV判据下的各分层响应数。每行分母为4，同一事件的0/1配对不增加独立事件数。}{p:counts}
\smallnote{主矩阵全部完成，最终无缺失响应。阈值改变主要在$\alpha$存储1的误读与临界状态之间重新分配；数值不确定分别为1、4、1项。}

\pageheading{测试前固定的预测}
训练集为$L=0.18/0.35\um$，$0.25\um$与$0.45\um$分别留作域内和外推测试。对两位的$\Delta Q$与$\Delta M$分别拟合，特征为$(Q_{\rm dep},Q_{\rm dep}/L,Q_{\rm dep}\sin\theta)$。

八个物理工况做留一交叉验证，每折留出完整工况的4个事件，标准化只使用该折训练数据。候选$\lambda=0,0.01,0.1,1,10$，按八折RMSE平均选择，四模型均取0.1；截距不惩罚。模型及64条预测在首次测试TCAD求解前固定，测试后未重训。
\begin{table}[H]\centering\small\caption{独立测试误差。每行$n=16$，无缺失或无效响应。}\scoretable\end{table}
\fig{prediction-residuals.png}{.90}{留出和外推残差的合并展示。分区数值见表，均值模型作为固定基准。}{p:prediction}
训练阶段四目标的岭回归交叉验证RMSE依次为0.102888 fC、0.955160 mV、0.112903 fC、1.012491 mV。独立测试的八项RMSE均小于相应均值基准；原评分与数值检查后的连续目标评分一致。

\pageheading{氧化层与数值细化}
固定$0.25\um$、垂直入射的一个$\alpha$和一个质子事件，对$4/6\,\mathrm{nm}$厚度重建网格及匹配参考；$5\,\mathrm{nm}$复用主结果。14次新计算得到8个事件响应，未增加独立事件数。
\fig{oxide-sensitivity.png}{1}{固定硅入口事件下的氧化层比较。细小差异采用下表数值解释。}{p:oxide}
\atable{oxide}{相对于5 nm结果的变化。}{p:oxidetable}
最大绝对变化为约0.005783687 fC和0.112560497 mV，分类保持。幅度小于0.01 fC和0.5 mV验收容差，但容差并非统计误差条；未对每个厚度追加独立细化和统计校准，不作微小效应显著性判断。

\atable{boundary}{两个边界目标的半步长和0.7网格比较。}{p:boundary}
\pageheading{读出边界与电场分布}
边界目标按完整128响应距三阈值的最小距离选取。四项连续比较均通过，最大电压差0.032654 mV、电荷差0.003490 fC；事件202610304的半步长结果在5 mV处跨界。它与3项未选近界响应合计为4项数值不确定。
全部384个响应--阈值对中，8项细化稳定、1项跨界、5项近界待确认、370项界外未细化，无连续量失败。
\fig{threshold-window-l045.png}{1}{$0.45\um$响应的3.8至4 ns裕量窗口。}{p:window}
\fig{native-event-window.png}{.96}{SVisual原生窗口：$0.18\um$、$\alpha$、$30^\circ$、事件202610311、存储1，$4\ns$电场幅值，单位V/cm。}{p:native}
粒子产生电荷、端口收集电荷与电容净变化受不同过程控制。电场图用于查看器件场分布，分类结论来自匹配参考下的BL差分及数值检查。

\pageheading{讨论与结论}
主阈值下，误读和未定状态集中于$\alpha$存储1。四个固定回归模型在独立尺寸测试中均低于均值基准误差，说明所选沉积、尺寸和方向特征对本组连续响应有预测作用。模型仍需真实Geant4沉积和无粒子基线，低残差也不能唯一确定微观电荷收集机制。

每层只有4个独立事件。分类统计保留固定分母，对临界和不确定项给出最佳/最差范围及Wilson 95\%区间；零次明确误读不足以支持零风险结论。数值细化中出现的跨界说明，连续量误差容差和阈值分类需要分开评价。

模型限制集中在二维平面结构、理想电容、准静态初态与总量等效径迹。所选质子电磁配置未验证核反应次级粒子效应，氧化层变化未重算完整材料栈输运。结果限于固定能量、入口、方向和栅长；实际写入、保持、感放锁存与恢复尚需器件数据校准，本文不据此估算商业DRAM失效率或持久位翻转截面。

本研究完成128个配对响应和64条固定测试预测的评价。后续实验应先校准器件与读出电路，再扩展能量、入射位置和独立样本。阈值附近仍应保留分类不确定，不能仅凭连续量差异较小将其清除。

\subsection*{计算记录与方法披露}
历史计算保留7条FAILED，其中主矩阵4次原始工程失败经相同输入重试后得到有效响应，原退出码未观测时仍为null。重试和18次扩展均不增加独立N。累计求解100022.07 s，包含失败和历史版本。

AI用于文献整理、脚本辅助、结果检查及文稿改写；指标与评分由固定程序计算。没有人工对照实验，也不据此报告AI加速率。逐事件结果、固定评分、失败记录和图件来源保存在随附数据\cite{project}。

\clearpage
\bibliographystyle{unsrt}
{\small\bibliography{../assets/references}}
\end{document}
